%-----------경력-----------%
% 불릿 작성 방식: 무엇을 만들었는지 먼저 쓰고, 그 결과를 숫자로 붙입니다.
% "결제 파이프라인을 이벤트 기반으로 재구축하여 응답 시간을 1.8초에서 0.4초로 단축"
% 처럼 쓰고, "Redis — 캐싱에 사용하여 성능 개선" 같은 나열식은 피합니다.
\section{경력 사항}
\resumeSubHeadingListStart

\resumeSubheading
  {Northwind Labs}{2023년 3월 -- 현재}
  {시니어 백엔드 엔지니어}{서울}
  \resumeItemListStart
    \resumeItem{결제 파이프라인을 NestJS와 BullMQ 기반의 이벤트 구조로 재구축하여 주문 처리 시간을 1.8초에서 0.4초로 단축했습니다.}
    \resumeItem{상품 카탈로그에 Redis 캐싱 계층을 설계하여 피크 시간대 PostgreSQL 읽기 부하를 60\% 줄였습니다.}
    \resumeItem{모놀리식 서비스를 6개의 배포 단위로 분리하는 마이그레이션을 주도하여 무중단으로 전환했습니다.}
    \resumeItem{주니어 엔지니어 3명을 주간 설계 리뷰로 멘토링했으며, 그중 2명이 같은 해에 승진했습니다.}
  \resumeItemListEnd

\resumeSubheading
  {BlueHarbor}{2022년 1월 -- 2023년 2월}
  {풀스택 개발자}{서울}
  \resumeItemListStart
    \resumeItem{React와 WebSocket으로 실시간 운영 대시보드를 구축하여 현재 물류 담당자 200명 이상이 매일 사용하고 있습니다.}
    \resumeItem{GitHub Actions로 빌드와 배포를 자동화하여 수작업 2시간이 걸리던 릴리스를 15분으로 단축했습니다.}
    \resumeItem{결제와 재고 모듈에 통합 테스트를 추가하여 백엔드 테스트 커버리지를 34\%에서 81\%로 끌어올렸습니다.}
    \resumeItem{REST API를 GraphQL로 전환하여 모바일 클라이언트의 중복 네트워크 호출을 40\% 제거했습니다.}
  \resumeItemListEnd

\resumeSubheading
  {Meridian Soft}{2021년 6월 -- 2021년 12월}
  {주니어 개발자}{부산}
  \resumeItemListStart
    \resumeItem{월 5,000명이 사용하는 고객 지원 포털을 와이어프레임 단계부터 배포까지 직접 맡아 출시했습니다.}
    \resumeItem{느린 MongoDB 집계 쿼리에 복합 인덱스를 적용하여 평균 페이지 로딩 시간을 3.2초에서 0.9초로 줄였습니다.}
  \resumeItemListEnd

\resumeSubHeadingListEnd
